\documentclass[11pt]{article}
\usepackage[margin=0.85in]{geometry}
\usepackage{amsmath,amssymb,booktabs,xurl,hyperref}
\usepackage[T1]{fontenc}
\hypersetup{colorlinks=true,urlcolor=blue}
\setlength{\emergencystretch}{2em}
\title{Why the Current CLF Stage Fails to Produce Nominal Recovery}
\author{Pan Dynamics Collision Avoidance Research}
\date{September 30, 2026}
\begin{document}
\maketitle

\section{Finding and scope}
The controller does execute its CLF stage. The principal failure isolated here
is that a useful first-input correction is almost erased while backtracking
the entire trajectory to satisfy a very small nonlinear endpoint core. The
remaining update is accepted because it slightly improves the incumbent's
relaxation, even though it increases the actual nominal-error energy. In
addition, the weighted secondary objective permits positive CLF relaxation
when a zero-relaxation trajectory exists. Straight terminal completion and a
shrinking horizon introduce this conflict even at exact circular cruise.

These statements are supported by new same-state interventions, rather than
inferred only from a long trajectory. Production source, configuration and
fixtures remain unchanged at commit
\path{17a5a19e76888079f8c7d0d0abc762a3b4150ca9}.
The probe has no target and no road constraints. It starts from exact trim on
a circle of radius 200 m at 8 m/s. The production controller executes 100
50-ms holds; tight ODE45 integration supplies its successors. Four captured
states are inspected without changing their original hard constraints.
There is no prescribed recovery deadline and no new online controller method.

\section{The CLF is locally capable of descent}
The nominal reference uses a discrete LQR matrix $P$ constructed from the
continuous transverse linearization and its exact zero-order-hold
transition. With feedback $K$ (including its negative sign), the Riccati
identity is
\[
P-(A+BK)^TP(A+BK)=Q+K^TRK.
\]
The smallest generalized eigenvalue of the left-hand side relative to $P$
is the guaranteed relative decrease of that local linear feedback model.
\begin{center}
\begin{tabular}{rrr}
\toprule
Speed (m/s)&Curvature (1/m)&Minimum local decrease per hold\\
\midrule
8&0&0.0409523\\
8&0.005&0.0409160\\
15&0&0.0336842\\
15&0.005&0.0326556\\
\bottomrule
\end{tabular}
\end{center}
The relative Riccati residuals are between $7.3\times10^{-16}$ and
$2.9\times10^{-15}$. All four local decreases exceed the configured
$\alpha=0.01$. Thus the required local decay is not unusually aggressive.
This is a local linear result, not a global nonlinear CLF proof. The stronger
relevant numerical evidence below is that the actual nonlinear model admits
zero-slack complete trajectories at all four inspected states.

\section{Where departure starts: the nominal state is not preserved}
The reference CLF uses the given-path curvature, whereas
\path{controller/collisionAvoidanceController.m} builds the terminal family
with curvature zero. Its pose is free; there is no fixed global endpoint
position or heading. However, the intrinsic terminal state is near straight
cruise: body velocities, yaw rate and last input are constrained near the
straight trim. At 8 m/s the curved nominal yaw rate is 0.04 rad/s and steering
is 0.0156894681 rad, rather than the straight trim's zero values.

The initial seed follows the nominal feedback for
$8+\lceil3/0.05\rceil=68$ holds and then appends straight-core completion.
It initially has 78 holds. Shifting consumes this plan; only when the count
would fall below eight is one terminal-feedback input appended. The initial
three-second recovery tail is not maintained as a moving preview.
At frame 69, time 3.40 s, the first straight-completion input reaches the
front of the shifted plan. The nominal energy before that hold is
$4.81\times10^{-27}$, effectively zero, but the issued steering is
$-0.0035098611$ rad and the next energy is 0.0462764. The horizon has ten
holds and reaches its minimum eight at frame 71.

This is a direct loss of nominal maintenance without an obstacle. It is not
proof that a straight terminal family necessarily forbids curved tracking:
complete zero-slack trajectories satisfying the same family are exhibited
below. Rather, the current initialization, horizon update and incomplete
correction turn a possible future straight continuation into an executed
maneuver that the CLF stage fails to undo.

\section{A same-frame causal intervention}
Consider frame 71 at time 3.50 s, with eight predicted holds. Its current
energy is $W=0.0678798354$; zero CLF slack requires
$W^+\le0.99W=0.0672010371$.
\begin{center}
\begin{tabular}{lrrl}
\toprule
Candidate&Next energy $W^+$&CLF slack&Original hard constraints\\
\midrule
Issued production candidate&0.104943902&0.037742865&Pass\\
Raw conic input sequence&0.052410637&0&Endpoint fails\\
Raw first input, corrected later inputs&0.052410637&0&Pass\\
LQR first input, corrected later inputs&0.027323953&0&Pass\\
\bottomrule
\end{tabular}
\end{center}
The raw conic sequence's only hard defect is its endpoint membership. Its
stage physical residual and safety slack are both zero. The terminal core
radius is $0.0001220703125$ in its weighted norm, not meters. The raw terminal
violation is 0.0286784 in that norm; its endpoint body-speed discrepancy is
only about $-0.003794$ m/s. This tight core is a sufficient local contraction
construction and must not be confused with the separate 6-mm collision
margin.

The controller checks scales $1,1/2,\ldots$ of the \emph{entire} input update.
At $1/4$ the next energy is about 0.03018, still a useful descent, but the
terminal norm fails. The first admitted scale is $1/128$, by which point
$W^+$ has risen to 0.1049439. Its slack improves from the shifted incumbent's
0.0398136 to 0.0377429, and that is sufficient to finish the online CLF stage.

For a discriminating intervention, the first raw control
$[0.0424190572,\ 0.0126863962]^T$ was held exactly fixed while only the later
seven controls were adjusted offline to reach the original endpoint center.
The resulting trajectory passes the unchanged nonlinear admission test.
Its largest input change is only 0.550 of the original input trust allowance;
its largest state change is 0.1025 of the original state trust allowance.
It therefore does not require relaxing the core or expanding those bounds.
The first hold independently replayed with ODE45 yields $W^+=0.0524128457$,
still well below the zero-slack threshold. Dense replay of its complete
horizon also remains inside the physical limits and final core.

The observed obstruction is therefore the numerical treatment of the
nonlinear endpoint, not unavoidable physical conflict or absence of a
sufficiently dissipative first input. The offline least-squares interpolation
is an existence and causality probe, not a proposed online nonlinear-shooting
method, backup controller or measured real-time implementation.

\section{Why repeating the existing iteration is insufficient}
At a feasible anchor, \path{controller/solvePredictiveControl.m} skips its
endpoint-polishing routine. The only response to the trial endpoint defect
is whole-trajectory backtracking. A linearized feasible direction generally
has a second-order nonlinear defect: for a smooth active endpoint row,
\[
g(U+d)=g(U)+Dg(U)d+O(\|d\|^2).
\]
A tiny terminal interior permits only a small uncorrected step. Returning
near the incumbent removes the endpoint defect by also removing the useful
first-input correction. The measured defect decreases roughly quadratically
with step scale, consistent with this mechanism.

At the same fixed states, thirty successive calls to the unchanged local
secondary operation give:
\begin{center}
\begin{tabular}{rrrrr}
\toprule
Frame&First admitted scale&Incumbent slack&After one update&After 30 updates\\
\midrule
69&1/32&0.0496066&0.0462764&0.0229515\\
71&1/128&0.0398136&0.0377429&0.0100348\\
80&1/64&0.2977394&0.2817460&0.2318087\\
100&1/128&2.3308371&2.3077506&2.2042249\\
\bottomrule
\end{tabular}
\end{center}
In each case the energy still increases after thirty updates. Nevertheless,
all four states admit complete zero-slack trajectories using an LQR first
input and adjusted later controls. At frame 100 that LQR first input exceeds
the single original input trust allowance; that distinction is recorded and
no single-step trust-feasibility claim is made there. Frames 71 and 80 already
supply stronger within-trust counterexamples.

The online code stops after one accepted secondary update. Its reported
\texttt{converged} and \texttt{clfStageCompleted} status means that numerical
stage contract was met, not that a nonlinear minimum or a Lyapunov decrease
was established. Increasing only the iteration cap leaves the demonstrated
small-step mechanism in place.

\section{A separate objective-level flaw}
Safety is optimized first, but CLF slack is not strictly prioritized over
tracking inside the secondary stage. The secondary criterion combines
\[
100\rho^2+\frac1H\sum_i W(e_i)
 +\frac{0.01}{H}\sum_i\|u_i-\bar u_i\|_R^2.
\]
The input penalty is correctly centered on the linearization inputs
$\bar u_i$, not zero. A finite squared-slack penalty does not require
$\rho=0$ whenever zero is feasible.

This can matter even at the exact nominal state. At frame 69 an offline
local optimization with the first input fixed at curved nominal trim obtains
objective 0.0388037973 and essentially zero slack. Freeing that first input
obtains a lower objective, 0.0386859983, with positive slack 0.0004378607.
Both trajectories pass the original hard constraints. Both offline solves
terminate on step tolerance; their stationarity residuals are about
$6\times10^{-7}$. These are local numerical results, not globally optimal
certificates. A separate fixed-steering profile also exhibits the trade:
a $+10^{-4}$ rad perturbation from nominal, with later controls reoptimized,
lowers the objective to 0.0387933372 while making slack positive.

Near $e=0$, a first-input perturbation of size $\epsilon$ gives
$W^+=O(\epsilon^2)$ and hence a squared CLF-slack penalty of
$O(\epsilon^4)$. That penalty need not dominate a first-order reduction in
future tracking/completion cost. Merely increasing a finite weight does not
supply nominal invariance or a strict CLF-priority theorem. Actual nonlinear
candidate comparison by slack does not change the objective used to generate
that candidate; lower-slack alternatives that were never generated cannot
be selected.

\section{The missing closed-loop condition and design consequences}
The executed relaxed inequality is
\[
W_{k+1}\le(1-\alpha)W_k+\rho_k.
\]
Improving a candidate's slack relative to another candidate at the
\emph{same state} says nothing about its size relative to $\alpha W_k$ or
about $\rho_k$ across successive states. Strict decrease requires
$\rho_k<\alpha W_k$. At frame 71, the issued slack is about 56 times
$\alpha W_k$, so its slight improvement cannot imply dissipation.

The last refactor corrected omission of the secondary solve, but did not
establish this closed-loop condition. A positive relaxation can be necessary
while avoiding a target. The missing property is eventual recovery of a
sufficient descent condition once compatible with safety, rather than merely
continuing to call the relaxed optimizer. No fixed fast settling time is
needed: a state-dependent positive-definite decrement is sufficient in an
appropriate invariant domain. A useful conditional alternative is
$\rho_k\to0$, since the geometric-convolution bound then yields $W_k\to0$.
The current stopping rule proves neither premise. Its no-solve numerical
threshold also admits a small nonzero relaxation, which by itself supports
at most a practical residual neighborhood.

The diagnosis points to three coupled design obligations within the same
controller: preserve useful first-input descent while restoring the original
nonlinear terminal constraints; define secondary-stage success in terms of
actual CLF priority/descent rather than any incumbent improvement; and make
the nominal continuation compatible with the reference path and maintained
prediction horizon. A zero-slack-first lexicographic objective can prevent
trading feasible CLF satisfaction for tracking cost, but still needs an
adequate nonlinear solver and an eventual-feasibility argument. An
affine-subproblem success alone is not enough. No backup design is required
by these findings.

The earlier extended campaign's ten straight-road cases did eventually meet
the sampled recovery criterion. Thus this report does not label every case
nonconvergent. The two indefinitely opposing same-circle head-on fixtures
also have recurring conflicts, a separate obstruction to permanent nominal
cruise. Neither fact explains away this target-free counterexample. The local
quadratic CLF is also not proved globally valid at the hundreds of meters of
error reached later; the present failure already appears near nominal trim,
before that global limitation becomes relevant.

\section{Validation and reproducibility}
Seven saved complete witnesses pass fresh unchanged nonlinear admission and
independent ODE45 replay at 31 samples per hold, including the final core and
physical state limits. The target-free safety-slack check contains no collision
rows. These are numerical sampled-model results, not physical-plant or
continuous-time safety certificates. No runtime speed claim is inferred from
the offline correction or extended 120-second diagnostic budget.

\path{report/CLF_CAUSAL_DIAGNOSIS_20260930/} contains compact summaries, the
line searches, all repeated-update outcomes, onset trace, objective profile,
Riccati checks, witness audit and exact raw/source hashes. Its reproduction
record lists the actual commands and cache-only instrumentation scripts.
Full snapshots and diagnostics remain under
\path{/home/zai/.cache/collisionAvoidance/clf-causal-diagnosis-20260930/}.
The three-panel figure is exported there as \texttt{causal\_diagnosis.pdf}
and \texttt{causal\_diagnosis.png}. No production code was changed and no
controller test-suite pass or full collision-campaign rerun is claimed for
this analysis-only task.
\end{document}
